% ============ 06. Vacíos críticos ============
\section{Los tres vacíos críticos (consenso de los tres evaluadores)}

\begin{enumerate}
  \item \textbf{S3 — no existe la declaración de uso de IA y reparto (½ pág.).}
        Vacío crítico confirmado por búsqueda exhaustiva. Sin ella (o con una
        falsa) S3 cae a Insuficiente. El borrador ya está redactado en el
        informe del agente S (\texttt{validacion\_anteproyecto\_S.md}, \S4.2)
        con tres casos verificables de ``lo que la IA sugirió y se descartó''.
  \item \textbf{No existe el anteproyecto de 6--10 págs.} Hay materia prima
        excelente (informe base, auditoría de 126 págs., SotA de 20 págs.),
        pero nada redactado en el formato exigido ni recontextualizado a
        2026-II / modalidad datos abiertos. Además hay una \textbf{ambigüedad
        conceptual D1/D3 sin decidir}: el objeto del anteproyecto debe fijarse
        (sistema/padrón de MinCiencias con datos abiertos — recomendado — vs.\
        el repositorio-código) antes de redactar problema y objetivos.
  \item \textbf{Preparación individual S1/S2 inexistente} (sin banco
        practicado, sin guion de 10 minutos, sin diapositivas del
        anteproyecto). Con el techo de autoría, es el mayor riesgo ponderado
        (la sustentación pesa 60\,\%).
\end{enumerate}

\section{Decisiones que debe tomar el equipo (no delegables)}

\begin{enumerate}
  \item \textbf{Objeto del anteproyecto} (D1/D3): se recomienda formularlo como
        \emph{``auditoría estadística reproducible del sistema de reconocimiento
        de investigadores de MinCiencias (2013--2021) con datos abiertos y
        propuesta de mejora del observatorio''}, en lugar de ``auditar un
        repositorio de código'' (que es el objeto de la consultoría ya
        ejecutada).
  \item \textbf{Modalidad: datos abiertos} — confirmar y asumir sus exigencias
        (pertinencia argumentada, procedencia con licencia y enlace, novedad
        frente a lo ya publicado con esos datos).
  \item \textbf{Reparto real del trabajo con nombres, roles y horas} (S3/D5) y
        \textbf{cronograma por hitos} de 12 semanas con 2 de holgura.
  \item \textbf{Número de integrantes} (afecta la frase final del reparto y el
        presupuesto de horas).
\end{enumerate}
